\chapter{Sprint 5 : Tableau de bord et IA}

\section{Introduction}
L'ultime sprint du projet se concentre sur l'expérience utilisateur finale, avec la réalisation des tableaux de bord interactifs et l'intégration profonde des fonctionnalités d'IA générative.

\section{Backlog du Sprint 5}
\begin{itemize}
    \item \textbf{US-13} : Développement des Dashboards (Admin, Teacher, Student).
    \item \textbf{US-14} : Intégration de la visioconférence (Daily.co).
    \item \textbf{US-15} : Finalisation des générateurs de résumés et de quiz IA.
\end{itemize}

\section{Analyse}
\subsection{Diagramme de cas d'utilisation}
L'Étudiant consulte son dashboard pour voir sa progression (XP). L'Enseignant lance une session vidéo. L'Administrateur surveille les statistiques globales.

\subsection{Description textuelle des cas d'utilisation}
L'enseignant clique sur "Générer Résumé". Le système appelle l'API Gemini, traite le résultat et l'affiche à l'écran. L'étudiant peut alors consulter ce résumé à côté de son support de cours.

\section{Conception}
\subsection{Diagrammes de séquence}
Le diagramme de séquence pour le lancement d'un Live montre l'interaction avec l'API Daily.co pour obtenir un token de salle.

\subsection{Diagramme de classes}
Classes \texttt{AIQuizController}, \texttt{AISummaryController} et les composants React du Dashboard.

\subsection{Représentation graphique de la base de données}
Documents stockant les résumés générés (\texttt{AISummary}) et les flashcards.

\section{Réalisation}
\subsection{Dashboard Étudiant avec Gamification}
Interface affichant les points d'expérience (XP), le GPA et les cours récents.
\begin{figure}[H]
    \centering
    \fbox{Placeholder: Capture d'écran - Student Dashboard}
    \caption{Tableau de bord intelligent de l'étudiant}
\end{figure}

\subsection{Interface de Visioconférence Intégrée}
Salle de classe virtuelle avec chat et liste des participants en temps réel.
\begin{figure}[H]
    \centering
    \fbox{Placeholder: Capture d'écran - Live Session}
    \caption{Module de visioconférence Daily.co}
\end{figure}

\section{Tests}
\subsection{Tests unitaires}
Tests des fonctions de formatage des réponses de l'IA.
\subsection{Tests d'intégration}
Validation de la persistance des résumés générés dans MongoDB.
\subsection{Tests fonctionnels}
Scénario complet : Upload d'un PDF $\rightarrow$ Génération d'un résumé $\rightarrow$ Consultation par l'étudiant $\rightarrow$ Passage du quiz lié.

\section{Conclusion}
Le Sprint 5 clôture la phase de développement avec une plateforme complète, intelligente et fonctionnelle, prête pour la mise en production.
